\documentclass[11pt,a4paper]{article}

\usepackage{fontspec}
\usepackage{polyglossia}
\setmainlanguage{russian}
\setotherlanguage{english}
\defaultfontfeatures{Ligatures=TeX}
\setmainfont{LinLibertine_R.otf}[
  BoldFont       = LinLibertine_RB.otf,
  ItalicFont     = LinLibertine_RI.otf,
  BoldItalicFont = LinLibertine_RBI.otf
]
\newfontfamily\cyrillicfonttt{LinLibertine_R.otf}[
  BoldFont       = LinLibertine_RB.otf,
  ItalicFont     = LinLibertine_RI.otf,
  BoldItalicFont = LinLibertine_RBI.otf
]

\usepackage[a4paper,top=2.2cm,bottom=2.2cm,left=2.3cm,right=2.3cm]{geometry}
\usepackage{graphicx}
\usepackage{booktabs}
\usepackage{amsmath}
\usepackage{amssymb}
\usepackage{caption}
\usepackage{hyperref}
\usepackage{csquotes}
\usepackage{enumitem}
\usepackage{setspace}
\usepackage[round,authoryear]{natbib}
\usepackage{titlesec}

\titlespacing*{\section}{0pt}{1.0ex plus 0.3ex minus 0.2ex}{0.4ex plus 0.2ex}
\titlespacing*{\subsection}{0pt}{0.6ex plus 0.2ex minus 0.1ex}{0.2ex plus 0.1ex}
\titleformat*{\section}{\normalfont\large\bfseries}
\titleformat*{\subsection}{\normalfont\normalsize\bfseries}

\setlist[itemize]{topsep=1pt, partopsep=0pt, itemsep=1pt, parsep=0pt, leftmargin=1.5em}
\setlist[enumerate]{topsep=1pt, partopsep=0pt, itemsep=1pt, parsep=0pt, leftmargin=1.7em}

\setlength{\abovecaptionskip}{2pt}
\setlength{\belowcaptionskip}{0pt}
\setlength{\textfloatsep}{4pt plus 2pt minus 2pt}
\setlength{\intextsep}{4pt plus 2pt minus 2pt}
\setlength{\floatsep}{4pt plus 2pt minus 2pt}

\hypersetup{
  colorlinks=true,
  linkcolor=black,
  citecolor=black,
  urlcolor=blue
}

\graphicspath{{outputs/proposal_eda/}}

\setlength{\parskip}{0.15em}
\setlength{\parindent}{1.25em}
\setlength{\emergencystretch}{3em}
\linespread{1.05}\selectfont

\title{Эмпирическая оценка влияния степени загрязнения воздуха на ожидаемую продолжительность жизни в регионах Российской Федерации}
\author{Ковалёва М.\,С., Кутузов И.\,В., Николюк Т.\,Н.}
\date{Москва, 2026}

\begin{document}

\begin{titlepage}
\centering
\vspace*{1cm}
{\large Федеральное государственное автономное образовательное учреждение\\
высшего образования\par}
\vspace{0.5cm}
{\Large \textbf{НАЦИОНАЛЬНЫЙ ИССЛЕДОВАТЕЛЬСКИЙ УНИВЕРСИТЕТ}\par}
\vspace{0.3cm}
{\Large \textbf{<<ВЫСШАЯ ШКОЛА ЭКОНОМИКИ>>}\par}
\vspace{0.5cm}
{\large Факультет экономических наук\par}
\vspace{0.3cm}
{\large Эконометрика\par}

\vfill

{\Large\bfseries Эмпирическая оценка влияния степени загрязнения воздуха на ожидаемую продолжительность жизни в регионах Российской Федерации\par}

\vfill

\begin{flushright}
\begin{minipage}{0.55\textwidth}
\textbf{Выполнили:}\\
Ковалёва Мария Сергеевна\\
Кутузов Иван Владимирович\\
Николюк Татьяна Николаевна\\[0.6em]
\textbf{Преподаватель семинаров:}\\
Бывальцева-Станкевич Анастасия Александровна
\end{minipage}
\end{flushright}

\vfill
{\large г.~Москва, 2026 г.\par}
\end{titlepage}

\section*{Введение}
\addcontentsline{toc}{section}{Введение}

Загрязнение атмосферного воздуха~--- одна из наиболее значимых экологических и социально-экономических проблем современности. По оценкам Всемирной организации здравоохранения~\cite{who2018}, загрязнение воздуха связано с повышенным риском сердечно-сосудистых, респираторных и онкологических заболеваний, что непосредственно негативно влияет на качество и ожидаемую продолжительность жизни. Для России данная проблема вызывает особый интерес вследствие высокой неоднородности регионов по уровню промышленной нагрузки, энергетической структуры экономики и концентрации вредных выбросов. В ряде субъектов РФ сохраняется высокий уровень загрязнения атмосферы, связанный с деятельностью металлургии, теплоэнергетики, добывающей промышленности и тяжёлого производства.

Таким образом, становится важным в рамках данного исследования определить влияние загрязнения воздуха на ожидаемую продолжительность жизни в регионах РФ для возможной корректировки политики в сторону улучшения экологической обстановки, а также повышения осведомлённости населения об угрозе для жизни в их регионе, представленной высокой концентрацией токсичных веществ в атмосферном воздухе.

\section*{Цель исследования}
\addcontentsline{toc}{section}{Цель исследования}

Цель настоящего эмпирического исследования заключается в оценке статистической значимости, силы и направления влияния степени загрязнения воздуха на ожидаемую продолжительность жизни населения по субъектам Российской Федерации в рамках эконометрической модели с учётом проблемы эндогенности.

Дополнительно исследование направлено на:
\begin{itemize}[leftmargin=2em]
  \item выявление факторов, связанных с региональными различиями в продолжительности жизни;
  \item проверку наличия эндогенности между загрязнением и социально-экономическими характеристиками регионов;
  \item оценку применимости инструментальных переменных для идентификации причинного эффекта загрязнения.
\end{itemize}

\section*{Задачи}
\addcontentsline{toc}{section}{Задачи}

Для достижения поставленной цели в работе предполагается решить следующие задачи:
\begin{enumerate}[leftmargin=2em]
  \item Изучить научную литературу по влиянию загрязнения воздуха на здоровье населения и продолжительность жизни.
  \item Сформировать и очистить кросс-секционный датасет по регионам Российской Федерации на основе данных Росстата и Fedstat.
  \item Провести предварительный анализ данных и исследовать основные статистические зависимости между переменными.
  \item Предложить эконометрическую спецификацию модели влияния загрязнения воздуха на ожидаемую продолжительность жизни.
  \item Определить потенциальные источники эндогенности и подобрать инструментальную переменную.
  \item Оценить модель методами OLS и 2SLS.
  \item Сравнить результаты различных спецификаций модели и провести проверки устойчивости результатов.
  \item Проинтерпретировать полученные оценки и сделать выводы о влиянии загрязнения воздуха на ожидаемую продолжительность жизни населения.
\end{enumerate}

\section*{Актуальность исследования}
\addcontentsline{toc}{section}{Актуальность исследования}

Проблема загрязнения воздуха остаётся актуальной как для мировой экономики, так и для России. По данным Всемирной организации здравоохранения, ухудшение качества воздуха связано с повышенным риском сердечно-сосудистых и респираторных заболеваний, а также с преждевременной смертностью.

Для России исследование данной темы особенно важно из-за существенной региональной неоднородности. Высокая концентрация добывающей и тяжёлой промышленности в отдельных регионах создаёт различия в экологической нагрузке, что потенциально может отражаться на состоянии здоровья населения.

Кроме того, тема представляет интерес с эконометрической точки зрения, поскольку загрязнение воздуха является потенциально эндогенной переменной: более индустриально развитые регионы одновременно могут иметь как более высокий уровень загрязнения, так и более высокий уровень доходов и качества медицины.

\section*{Гипотезы исследования}
\addcontentsline{toc}{section}{Гипотезы исследования}

\textbf{Гипотеза 1.} Рост выбросов загрязняющих веществ на душу населения сокращает ожидаемую продолжительность жизни в регионах РФ.
\begin{align*}
H_0\colon\ & \beta_1 = 0 \quad \text{(выбросы не влияют на ОПЖ)}\\
H_1\colon\ & \beta_1 < 0 \quad \text{(рост выбросов сокращает ОПЖ)}
\end{align*}
Ожидаемое направление: $\beta_1 < 0$ в спецификации (\ref{eq:ols}).

В обоснование гипотезы можно привести несколько работ. \citet{pope2002} показали на данных American Cancer Society, что рост долгосрочной концентрации $\text{PM}_{2.5}$ на 10~мкг/м³ увеличивает сердечно-сосудистую и лёгочную смертность примерно на 6\% и 8\%. \citet{chen2013} и \citet{ebenstein2017} на данных о реке Хуайхэ оценили потери продолжительности жизни на севере Китая в 3--5 лет, задавая тем самым верхнюю границу для российского контекста с его более умеренным уровнем загрязнения. \citet{currie2005} добавляют, что значимые эффекты на здоровье фиксируются даже при относительно невысоких уровнях экспозиции (низких уровнях загрязнения воздуха). Наконец, \citet{revich1998} оценили избыточную смертность от атмосферных выбросов непосредственно для крупных городов России.

\textbf{Гипотеза 2.} После контроля на загрязнение и возрастную структуру доход положительно связан с ОПЖ.
\begin{align*}
H_0\colon\ & \beta_{\text{grp\_pc}} = 0 \quad \text{(доход не связан с ОПЖ при прочих равных)}\\
H_1\colon\ & \beta_{\text{grp\_pc}} > 0 \quad \text{(более богатые регионы живут дольше при фиксированном уровне загрязнения)}
\end{align*}

\citet{preston1975} показал устойчивую положительную связь между национальным доходом на душу населения и ОПЖ в межстрановом разрезе, однако подчеркнул, что значительная часть вариации определяется факторами, экзогенными доходу. В российских регионах сырая корреляция отрицательна ($r = -0{,}13$), что отражает специфику сырьевых территорий: высокий ВРП здесь сопряжён с промышленной нагрузкой, суровым климатом и специфической структурой занятости, а не с теми механизмами улучшения здоровья, которые описывает Preston. Это подчёркивает необходимость многомерной спецификации с явным контролем дохода (отдельный регрессор).

\textbf{Гипотеза 3.} Занятость в добыче полезных ископаемых на душу населения связана с уровнем загрязнения воздуха после учёта контрольных переменных.
\begin{align*}
H_0\colon\ & \pi_1 = 0\\
H_1\colon\ & \pi_1 > 0
\end{align*}

\section*{Данные}
\addcontentsline{toc}{section}{Данные}

В работе были использованы кросс-секционные данные по 85 регионам России за 2015 год, так как именно за этот год данные имелись в наиболее полном объёме. Основной источник данных~--- Fedstat, ссылки на конкретные датасеты представлены в приложении. Все переменные, описанные ниже, являются непрерывными количественными. Также используется переменная \texttt{region}~--- категориальная. В моделях используются логарифмы некоторых переменных.

\textbf{Зависимая переменная.} \texttt{life\_expectancy}~--- ожидаемая продолжительность жизни в годах, основной показатель здоровья населения.

\textbf{Основная объясняющая переменная.} \texttt{emissions\_h1\_pc}~--- выбросы загрязняющих веществ от стационарных источников на душу населения, тыс.~тонн на человека.

\textbf{Инструментальные переменные:}
\begin{itemize}
  \item \texttt{emp\_mining\_pc}~--- занятость в добыче полезных ископаемых на душу населения, основной инструмент;
  \item \texttt{emissions\_fuel\_pc}~--- выбросы от сжигания топлива на душу населения;
  \item \texttt{emp\_energy\_pc}~--- занятость в энергетике на душу населения.
\end{itemize}

\textbf{Контрольные переменные:}
\begin{itemize}
  \item \texttt{grp\_pc}~--- ВРП на душу населения, тыс.~руб. на человека;
  \item \texttt{urban\_share}~--- доля городского населения, \%;
  \item \texttt{age\_struct\_share}~--- доля населения старших возрастов, \%;
  \item \texttt{population\_avg}~--- среднегодовая численность населения, человек.
\end{itemize}

\textbf{Дополнительные переменные:}
\begin{itemize}
  \item \texttt{industrial\_index}~--- индекс промышленного производства;
  \item \texttt{visits\_per\_capita}~--- число посещений лечебных учреждений на одного жителя;
  \item \texttt{emp\_manufacturing\_pc}~--- занятость в обрабатывающих производствах на душу населения;
  \item \texttt{emp\_construction\_pc}~--- занятость в строительстве на душу населения;
  \item \texttt{emp\_metallurgy\_pc}~--- занятость в металлургии на душу населения;
  \item \texttt{emp\_chemicals\_pc}~--- занятость в химическом производстве на душу населения.
\end{itemize}

\section*{Методы и спецификация модели}
\addcontentsline{toc}{section}{Методы и спецификация модели}

Оценивание проводится в два шага. Сначала строится базовая OLS-регрессия~--- она даёт состоятельные оценки только при отсутствии эндогенности, которая здесь может присутствовать. Затем применяется 2SLS, где в качестве инструмента для эндогенной переменной \texttt{ln\_emissions\_h1\_pc} выступает логарифм занятости в добыче полезных ископаемых на душу населения. Во всех спецификациях стандартные ошибки скорректированы на гетероскедастичность (HC1).

Базовая OLS-регрессия (уравнение 1):
\begin{equation}
\label{eq:ols}
\begin{aligned}
\text{life\_expectancy}_i = {} & \beta_0 + \beta_1\,\text{ln\_emissions\_h1\_pc}_i + \beta_2\,\text{ln\_grp\_pc}_i \\
& + \beta_3\,\text{urban\_share}_i + \beta_4\,\text{age\_struct\_share}_i \\
& + \beta_5\,\text{ln\_population\_avg}_i + \varepsilon_i
\end{aligned}
\end{equation}

Первая ступень 2SLS (уравнение 2):
\begin{equation}
\label{eq:fs}
\begin{aligned}
\text{ln\_emissions\_h1\_pc}_i = {} & \pi_0 + \pi_1\,\text{ln\_emp\_mining\_pc}_i + \pi_2\,\text{ln\_grp\_pc}_i \\
& + \pi_3\,\text{urban\_share}_i + \pi_4\,\text{age\_struct\_share}_i \\
& + \pi_5\,\text{ln\_population\_avg}_i + v_i
\end{aligned}
\end{equation}

Вторая ступень 2SLS (уравнение 3):
\begin{equation}
\label{eq:ss}
\begin{aligned}
\text{life\_expectancy}_i = {} & \beta_0 + \beta_1\,\widehat{\text{ln\_emissions\_h1\_pc}}_i + \beta_2\,\text{ln\_grp\_pc}_i \\
& + \beta_3\,\text{urban\_share}_i + \beta_4\,\text{age\_struct\_share}_i \\
& + \beta_5\,\text{ln\_population\_avg}_i + \varepsilon_i
\end{aligned}
\end{equation}

Логарифмическая форма для выбросов, ВРП и численности населения выбрана из-за правосторонней асимметрии этих переменных. Коэффициент $\beta_1$ интерпретируется как полу-эластичность: изменение ОПЖ в годах при удвоении выбросов на душу населения равно $\beta_1\cdot\ln 2 \approx \beta_1\cdot 0{,}693$.

\section*{Эндогенность и инструментальная стратегия}
\addcontentsline{toc}{section}{Эндогенность и инструментальная стратегия}

OLS-оценка $\beta_1$ из уравнения~(\ref{eq:ols}) смещена как минимум в силу:
\begin{itemize}[leftmargin=2em]
  \item пропущенных переменных~--- качество здравоохранения, климат, исторически сложившаяся отраслевая структура связаны одновременно и с загрязнением, и с продолжительностью жизни;
  \item ошибок измерения в \texttt{emissions\_h1}, которая не учитывает транспорт, бытовое отопление и трансграничный перенос;
  \item миграционных причин~--- здоровые трудоспособные граждане концентрируются в индустриальных регионах, искусственно завышая там наблюдаемую продолжительность жизни.
\end{itemize}

В качестве инструмента используется логарифм занятости в добыче полезных ископаемых на душу населения. Его релевантность объясняется тем, что размещение добывающих производств определяется геологией и историей освоения территорий, а не текущей политикой здравоохранения, тогда как добыча и первичная переработка сырья~--- один из главных источников стационарных выбросов. Ожидается, что $F$-статистика первой ступени превысит порог \citet{staiger1997} $F > 10$; в случае умеренной силы инструмента будут применены робастные к слабым инструментам методы.

\section*{Диагностика и робастность}
\addcontentsline{toc}{section}{Диагностика и робастность}

Планируется: (1) тест Wu--Hausman / Дурбина на экзогенность регрессора; (2) проверка силы инструмента по частной $F$-статистике с критическими значениями \citet{stockyogo2005}; (3) VIF-анализ мультиколлинеарности; (4) тест Саргана при использовании более одного инструмента; (5) переоценка без влиятельных наблюдений (Cook's distance, leverage); (6) альтернативные инструменты~--- занятость в энергетике.

\section*{Ограничения}
\addcontentsline{toc}{section}{Ограничения}

Кросс-секция за один год не позволяет контролировать ненаблюдаемую региональную гетерогенность через фиксированные эффекты~--- это недостаток, который лишь частично компенсирует инструментальная переменная. Данные об \texttt{emissions\_h1} охватывают только стационарные источники, исключая транспорт и бытовой сектор, что занижает суммарную экспозицию.

\section*{Предварительный анализ данных}
\addcontentsline{toc}{section}{Предварительный анализ данных}

\begin{table}[ht]
\centering
\caption{Дескриптивные статистики основных переменных}
\label{tab:desc}
\begin{tabular}{lrrrrr}
\toprule
\textbf{Переменная} & \textbf{Среднее} & \textbf{Ст.\,откл.} & \textbf{Мин} & \textbf{Медиана} & \textbf{Макс} \\
\midrule
\texttt{life\_expectancy}    & 70{,}5662  & 2{,}3977  & 63{,}1400   & 70{,}3900  & 78{,}4700 \\
\texttt{ln\_emissions\_h1\_pc} & $-10{,}3097$ & 1{,}4159 & $-14{,}5157$ & $-10{,}2621$ & $-6{,}9957$ \\
\texttt{ln\_emp\_mining\_pc}  & $-5{,}5818$  & 1{,}6311 & $-8{,}7341$  & $-5{,}9033$  & $-1{,}5898$ \\
\texttt{ln\_grp\_pc}          & 5{,}8500   & 0{,}6917  & 4{,}6834    & 5{,}7689   & 8{,}5867 \\
\texttt{urban\_share}        & 70{,}2682  & 13{,}2684 & 29{,}9000   & 71{,}4000  & 100{,}0000 \\
\texttt{age\_struct\_share}  & 58{,}2165  & 2{,}0328  & 54{,}2000   & 57{,}7000  & 64{,}9000 \\
\texttt{ln\_population\_avg} & 13{,}9713  & 0{,}9782  & 10{,}6546   & 13{,}9885  & 16{,}3233 \\
\bottomrule
\end{tabular}
\end{table}

\begin{figure}[htbp]
\centering
\begin{minipage}[t]{0.49\textwidth}
\centering
\includegraphics[width=\linewidth]{01_life_expectancy_distribution.png}
\caption{Распределение ожидаемой продолжительности жизни по регионам РФ}
\label{fig:le}
\end{minipage}\hfill
\begin{minipage}[t]{0.49\textwidth}
\centering
\includegraphics[width=\linewidth]{02_emissions_distribution.png}
\caption{Распределение логарифма выбросов загрязняющих веществ на душу населения}
\label{fig:em}
\end{minipage}
\end{figure}

\begin{figure}[htbp]
\centering
\begin{minipage}[t]{0.49\textwidth}
\centering
\includegraphics[width=\linewidth]{03_life_expectancy_vs_emissions.png}
\caption{Связь между загрязнением воздуха и ожидаемой продолжительностью жизни}
\label{fig:le_em}
\end{minipage}\hfill
\begin{minipage}[t]{0.49\textwidth}
\centering
\includegraphics[width=\linewidth]{04_emissions_vs_mining_instrument.png}
\caption{Связь между занятостью в добыче полезных ископаемых и загрязнением воздуха}
\label{fig:em_mining}
\end{minipage}
\end{figure}

\begin{figure}[htbp]
\centering
\includegraphics[width=0.55\textwidth]{05_correlation_heatmap.png}
\caption{Корреляционная матрица основных переменных}
\label{fig:corr}
\end{figure}

Предварительный анализ данных позволил сделать несколько содержательных выводов: по некоторым показателям, таким как загрязнение воздуха, значения существенно различаются, что объясняется тем, что в регионах разный уровень развития промышленности. Распределение выбросов заметно асимметрично, так что данная переменная обоснованно прологарифмирована. Между ОПЖ и загрязнением воздуха наблюдается умеренная отрицательная корреляция. Занятость в добыче положительно связана с загрязнением, что подтверждает релевантность инструмента.

\nocite{*}
\bibliographystyle{plainnat}
\bibliography{references}

\section*{Примечание}

При подготовке работы использовался Claude Code для преобразования данных, построения EDA и анализа результатов.

\end{document}
